\documentclass[11pt]{article}
\usepackage[a4paper,margin=2.3cm]{geometry}
\usepackage{amsmath,amssymb,bm}
\usepackage{booktabs,longtable,array}
\usepackage{enumitem}
\usepackage[dvipsnames]{xcolor}
\usepackage[most]{tcolorbox}
\usepackage{microtype}
\usepackage[colorlinks=true,linkcolor=NavyBlue,citecolor=ForestGreen,urlcolor=NavyBlue]{hyperref}
\setlist{itemsep=2pt,topsep=3pt}
\newtcolorbox{keybox}[1]{colback=NavyBlue!5,colframe=NavyBlue!70,fonttitle=\bfseries,title={#1},left=4pt,right=4pt,top=3pt,bottom=3pt}
\newtcolorbox{warnbox}[1]{colback=BrickRed!5,colframe=BrickRed!70,fonttitle=\bfseries,title={#1},left=4pt,right=4pt,top=3pt,bottom=3pt}
\newtcolorbox{ideabox}[1]{colback=ForestGreen!5,colframe=ForestGreen!70,fonttitle=\bfseries,title={#1},left=4pt,right=4pt,top=3pt,bottom=3pt}
\newcommand{\bb}{\mathbf{b}}
\newcommand{\h}{\mathbf{h}}
\title{\textbf{Sparse Arrays for Near-Field Cell-Free XL-MIMO}\\[2pt]
\large A plain-language guide, literature map, and plan for extending the ICASSP and VTC papers into a Transactions paper}
\author{Prepared for Nayyab Haider (PUC-Rio)}
\date{October 2026}

\begin{document}
\maketitle

\begin{keybox}{How to read this document}
Sections~\ref{sec:what}--\ref{sec:why} explain what a sparse array is and why it matters in the near field, using numbers computed for \emph{your} system ($L=8$ APs, $N=64$ antennas, $f_c=3$\,GHz, $300\times300$\,m$^2$).
Section~\ref{sec:rules} lists the rules you must respect. Section~\ref{sec:math} collects the equations.
Section~\ref{sec:lit} maps the literature, and Section~\ref{sec:gap} states the gap.
Sections~\ref{sec:options}--\ref{sec:plan} give concrete options and a step-by-step plan for the Transactions paper.
\end{keybox}

\begin{warnbox}{Scope of the literature search (read this first)}
The search was done through a web search engine. Direct access to arXiv, Google Scholar, IEEE Xplore, ResearchGate and the GitHub API was blocked by the network policy of the computing environment, so full texts could not be opened and GitHub code could not be browsed.
Every reference below was checked for title and venue; where the author list could not be verified it is marked ``authors not verified''. Before citing any of them in a paper, open the PDF and confirm the details and the exact claims.
``No paper found'' in Section~\ref{sec:gap} means ``not found by this search'', not ``does not exist''.
\end{warnbox}

\tableofcontents
\newpage

%=====================================================================
\section{Where your work stands, and what a Transactions paper needs}
\label{sec:context}

\textbf{ICASSP paper.} Cross-AP list-based SIC for hybrid near-/far-field cell-free XL-MIMO, with \emph{perfect} CSI. The list is formed from both the cluster-fused and the all-AP fused soft estimates, and cancellation is network-wide. An LDPC-coded IDD scheme is added.

\textbf{VTC paper.} The same detector, plus \emph{estimated} CSI from polar-domain channel estimation over all AP--UE pairs [1].

\textbf{What a Transactions paper (IEEE TWC or IEEE TCOM) needs.}
A journal reviewer will expect (i) a new technical dimension, not only a merge of the two conference papers, (ii) analysis, not only simulations, and (iii) a broad numerical study.
A sparse array at each AP is a strong candidate for this new dimension. It changes three things your papers depend on:
\begin{enumerate}
  \item \emph{where} the near field is (the Rayleigh distance grows with the square of the spacing);
  \item \emph{how well} users at the same angle but different range can be separated (this is the ``range-domain resolution'' your papers rely on);
  \item \emph{what goes wrong} (grating lobes), and the cross-AP idea of your detector is a natural way to fix it.
\end{enumerate}

%=====================================================================
\section{What is a sparse array?}
\label{sec:what}

A conventional uniform linear array (ULA) places its $N$ antennas at half-wavelength spacing $d=\lambda/2$. A \textbf{sparse array} uses spacing larger than $\lambda/2$, or places the antennas at non-uniform positions.

\paragraph{Main families.}
\begin{itemize}
  \item \textbf{Uniform sparse ULA}: spacing $d=\eta\,\lambda/2$ with sparsity factor $\eta>1$ (for example $\eta=2$ or $4$). Element positions
  \begin{equation}
    p_n=\Big(n-\tfrac{N+1}{2}\Big)\,\eta\,\frac{\lambda}{2},\qquad n=1,\dots,N.
  \end{equation}
  \item \textbf{Thinned array}: start from a dense ULA and switch off (or remove) some elements, so the aperture stays the same but fewer antennas and RF chains are used [10].
  \item \textbf{Non-uniform sparse arrays}: the minimum-redundancy array (MRA) [18], the nested array [19] and the co-prime array [20]. Their positions are chosen so that the pairwise differences $p_m-p_n$ cover many values. This gives $\mathcal{O}(N^2)$ virtual elements in the difference co-array, which is very useful for sensing [9].
  \item \textbf{Modular (or quasi-distributed) arrays}: several dense sub-arrays separated by large gaps [11], [12], [13], [14]. This is a sparse array at the level of sub-arrays.
\end{itemize}

\paragraph{The two ways to compare a sparse array with a dense array.}
These give very different conclusions, so always state which one you use:
\begin{itemize}
  \item \textbf{Same number of antennas $N$}: the sparse array has a larger aperture $D=(N-1)\eta\lambda/2$. You gain resolution and a larger near-field region, and you pay with grating lobes.
  \item \textbf{Same aperture $D$}: the sparse array uses about $N/\eta$ antennas. Resolution is (almost) unchanged, and you save hardware, power and computation. You still pay with grating lobes.
\end{itemize}

%=====================================================================
\section{Why sparse arrays matter in the near field: numbers for your system}
\label{sec:why}

\subsection{The near-field region grows with $\eta^2$}
The Rayleigh distance used in your papers (eq.~(2) of the VTC paper) is $R=2D^2/\lambda$. With $D=(N-1)\eta\lambda/2$,
\begin{equation}
  R(\eta)=\frac{(N-1)^2\eta^2\lambda}{2}=\eta^2\,R(1).
  \label{eq:rayleigh}
\end{equation}
\textbf{In plain words: doubling the spacing makes the near-field region four times larger.}

\subsection{Distance (range) focusing needs $r$ well inside $R$}
Being inside $R$ is not enough to get a sharp focus in distance. The near-field beamforming primer [3] shows that a beam has a \emph{finite depth} only for distances below about $R/10$; beyond that, the beam extends to infinity in range, just like a far-field beam.
For a focus at distance $F<R/\alpha$, the 3-dB depth interval has the form [3], [4]
\begin{equation}
  \frac{R\,F}{R+\alpha F}\;\le\; r\;\le\;\frac{R\,F}{R-\alpha F},
  \qquad
  \text{depth}=\frac{2\alpha F^2 R}{R^2-\alpha^2F^2},
  \label{eq:depth}
\end{equation}
where $\alpha\approx 10$ in [3]. The exact constant depends on the array geometry (values around 7 have been reported for ULAs), so treat $\alpha$ as ``between about 7 and 10''.

\begin{table}[h]
\centering
\caption{Numbers for your system ($N=64$ unless stated, $f_c=3$\,GHz, $\lambda=0.1$\,m). ``Corr.'' is the correlation $|\bb_1^{\mathsf H}\bb_2|$ between the unit-norm near-field array vectors of two users at the \emph{same} angle and different ranges. Small values mean the array can separate them by range. Computed with the exact NUSW model.}
\label{tab:numbers}
\begin{tabular}{@{}cccccccc@{}}
\toprule
$\eta$ & $N$ & spacing & aperture $D$ & $R$ & $R/10$ & Corr.\ 30\,m vs 120\,m & Corr.\ 40\,m vs 80\,m\\
 & & [m] & [m] & [m] & [m] & ($\sin\theta=0$) & ($\sin\theta=0$)\\
\midrule
1 & 64 & 0.05 & 3.15 & 198 & 19.8 & 0.83 & 0.96\\
2 & 64 & 0.10 & 6.30 & 794 & 79.4 & 0.36 & 0.46\\
4 & 64 & 0.20 & 12.6 & 3175 & 318 & 0.15 & 0.24\\
\midrule
2 & 32 & 0.10 & 3.10 & 192 & 19.2 & 0.83 & 0.96\\
4 & 16 & 0.20 & 3.00 & 180 & 18.0 & 0.83 & 0.96\\
\bottomrule
\end{tabular}
\end{table}

\begin{keybox}{What Table~\ref{tab:numbers} tells you about your current papers}
\begin{itemize}
  \item In your simulations each user is placed 20--150\,m from its anchor AP (other APs are farther). With compact APs ($\eta=1$), $R/10\approx 20$\,m, so \textbf{almost no link has a finite beam depth}. The links are ``near-field'' by the Rayleigh criterion, but the range-domain separation is weak: two co-angular users at 30\,m and 120\,m still have correlation 0.83.
  \item With the \emph{same 64 antennas} at $\eta=2$ the correlation drops to 0.36, and at $\eta=4$ to 0.15. The range-domain resolution that your papers talk about becomes \emph{real}.
  \item With the \emph{same aperture} (rows 4--5), the correlation does not change: range resolution depends on the aperture, not on the number of antennas. A thinned array gives the same separation with $2\times$ or $4\times$ fewer antennas.
\end{itemize}
\end{keybox}

\subsection{More degrees of freedom}
For a single-user MIMO link, the effective degrees of freedom (EDoF) grow with the array sparsity until they reach the maximum $\min(N_t,N_r)$ [8]. For multi-user detection, the analogous effect is that the local channel matrix $\mathbf{H}_l$ becomes better conditioned, which is exactly the ``rank enhancement'' of eq.~(10) in your papers.

\subsection{Grating lobes: the price}
A ULA with spacing $d>\lambda/2$ has an array factor that repeats in $\sin\theta$ with period $\lambda/d=2/\eta$. In the \emph{far field}, every repetition (a \emph{grating lobe}) has the \emph{same} gain as the main lobe. A user located at a grating-lobe angle cannot be separated from the intended user by that array.

In the \emph{near field}, the picture changes (Section~\ref{sec:gl-math}): the grating lobe is not removed, but it is \textbf{focused at a different range}. Example computed for your array at $\eta=2$: focusing at $(\sin\theta,r)=(0.3,40\,\text{m})$ gives
\begin{itemize}
  \item gain 0.38 at the grating-lobe angle $\sin\theta=-0.7$ \emph{at the same range} (suppressed);
  \item gain 0.89 at the grating-lobe angle at the \emph{shifted range} $r'=22.4$\,m (almost full);
  \item gain 1.00 for a \emph{far-field} user at the grating-lobe angle (no suppression at all).
\end{itemize}
This is consistent with the literature [5], [7], [11]: near-field operation suppresses grating lobes ``at the same distance'', but not everywhere.

\begin{ideabox}{Why this is good news for a cell-free paper}
In a hybrid-field cell-free network, the grating-lobe ``danger zone'' of a user is different at every AP, because each AP sees the user at a different angle and range. Two users that collide on a grating lobe at one AP are almost never colliding at all APs. A detector that fuses \emph{all} APs, as your cross-AP list SIC does, can resolve what a single AP (or the serving cluster) cannot. Sparse APs therefore make the \emph{information surplus} $\Delta_k$ of your papers (eq.~(26) in VTC, eq.~(22) in ICASSP) larger and more important.
\end{ideabox}

%=====================================================================
\section{Rules and conditions to respect}
\label{sec:rules}

\begin{enumerate}[label=\textbf{R\arabic*.}]
  \item \textbf{State the comparison basis.} Same $N$ or same aperture (Section~\ref{sec:what}). Reviewers reject comparisons where the sparse array silently gets a larger aperture \emph{and} the dense array is not given the same.
  \item \textbf{Recompute the near/far-field boundary.} $R$ scales with $\eta^2$ (eq.~\eqref{eq:rayleigh}). In your $300\times300$\,m$^2$ area, $\eta=2$ already makes almost every link near-field ($R=794$\,m), so the \emph{hybrid} regime disappears. To keep a hybrid network you need a larger area, mixed APs (some dense, some sparse), or a boundary that also uses the angle (the effective Rayleigh distance depends on $1-\sin^2\theta$ [2], [31]).
  \item \textbf{Range focusing needs $r \lesssim R/\alpha$}, $\alpha\approx 7$--$10$ (eq.~\eqref{eq:depth}). Check where your users are relative to this limit; report the fraction of links inside it.
  \item \textbf{Account for grating lobes in every model and algorithm}: in the inter-user correlation, in the channel estimator's dictionary, and in the clustering rule. Far-field links of a sparse AP see grating lobes at full strength.
  \item \textbf{Dictionary design for estimation.} The angular grid must be finer by $\eta$ (beams are $\eta$ times narrower), and atoms that differ by $2/\eta$ in $\sin\theta$ are nearly identical except for their curvature. Polar-domain OMP will confuse them unless the estimator explicitly tests the aliases (Section~\ref{sec:est}).
  \item \textbf{Keep the element model consistent.} In your code the far-field correlation matrix uses the spacing in wavelengths (\texttt{d\_H = 0.5}); it must become $\eta/2$ together with the NUSW positions. A mismatch would silently compare two different arrays.
  \item \textbf{Physical limits.} Large spacing reduces mutual coupling (good) but increases the physical size; element directivity reduces grating lobes near end-fire; wideband operation makes grating lobes split in frequency [9]. Your system is narrowband, so state this assumption.
  \item \textbf{Non-uniform arrays} (nested, co-prime, MRA) break the periodicity, so there are no exact grating lobes, but the side-lobes rise. They are attractive for estimation (larger co-array) [9], [17], [19], [20].
\end{enumerate}

%=====================================================================
\section{Mathematical background}
\label{sec:math}

\subsection{Channel model with a sparse AP}
Keep your NUSW model (eq.~(6) of the papers) and only change the positions:
\begin{equation}
  [\h^{\rm NF}_{l,k}]_n=\sqrt{\beta_{l,k}}\,\frac{d_{l,k}}{d_{n,l,k}}\,
  e^{-j\frac{2\pi}{\lambda}d_{n,l,k}},\qquad
  d_{n,l,k}=\sqrt{d_{l,k}^2+p_n^2-2\,d_{l,k}\,p_n\sin\theta_{l,k}},
\end{equation}
with $p_n$ from (1). The far-field model becomes $\sqrt{\beta_{l,k}}\,\mathbf a(\theta_{l,k})$ with $[\mathbf a(\theta)]_n=e^{j\frac{2\pi}{\lambda}p_n\sin\theta}$.

\subsection{Second-order (Fresnel) form and the array gain}
With $d_{n}\approx r-p_n s+\frac{p_n^2(1-s^2)}{2r}$, $s=\sin\theta$, the unit-norm near-field steering vector is
\begin{equation}
  [\bb(s,r)]_n\approx\frac{1}{\sqrt N}\exp\!\Big(j\frac{2\pi}{\lambda}p_n s-j\frac{\pi}{\lambda}\frac{p_n^2(1-s^2)}{r}\Big),
\end{equation}
and the normalized gain when focusing at $(s_0,r_0)$ and looking at $(s,r)$ is
\begin{equation}
  G(s,r)=\Big|\frac1N\sum_{n=1}^N
  \exp\!\Big(j\frac{2\pi}{\lambda}p_n(s-s_0)
  -j\frac{\pi}{\lambda}p_n^2\Big[\frac{1-s^2}{r}-\frac{1-s_0^2}{r_0}\Big]\Big)\Big|^2 .
  \label{eq:gain}
\end{equation}

\subsection{Grating lobes in the near field (simple derivation)}
\label{sec:gl-math}
For a uniform sparse ULA, $p_n=(n-\frac{N+1}{2})\eta\lambda/2$, so the linear phase in \eqref{eq:gain} is $\pi\eta\,(n-\frac{N+1}{2})(s-s_0)$. It is a multiple of $2\pi$ for every $n$ when
\begin{equation}
  s_m=s_0+\frac{2m}{\eta},\qquad m=\pm1,\pm2,\dots,\quad |s_m|\le 1 .
\end{equation}
At these angles only the quadratic term in \eqref{eq:gain} remains:
\begin{equation}
  \Phi_n=\frac{\pi p_n^2}{\lambda}\Big[\frac{1-s_0^2}{r_0}-\frac{1-s_m^2}{r}\Big].
\end{equation}
Two consequences follow, and both are visible in the numerical example of Section~\ref{sec:why}:
\begin{enumerate}
  \item \textbf{At the same range} ($r=r_0$), $\Phi_n=\pi p_n^2(s_m^2-s_0^2)/(\lambda r_0)$. Its largest value, at the array edge $p=D/2$, is
  \begin{equation}
    \Phi_{\max}=\frac{\pi R}{8\,r_0}\,\big|s_m^2-s_0^2\big| .
  \end{equation}
  The grating lobe is strongly attenuated when $\Phi_{\max}\gtrsim\pi$, that is when $r_0\lesssim R\,|s_m^2-s_0^2|/8$. Since $R\propto\eta^2$, sparser arrays and closer users give stronger suppression, which matches [5].
  \item \textbf{At the shifted range} $r'=r_0\,\dfrac{1-s_m^2}{1-s_0^2}$, $\Phi_n=0$ and the grating lobe is at \emph{full} gain (up to higher-order terms). This is an \emph{ambiguity point}: a second user located at $(s_m,r')$ cannot be separated by this AP.
\end{enumerate}
For a far-field user ($r\to\infty$) and a far-field focus, $\Phi_n\to0$: the grating lobe is always at full gain.

\subsection{Inter-user correlation and resolvability}
At AP $l$ define the correlation between users $i$ and $j$:
\begin{equation}
  \mu_l(i,j)=\frac{|\h_{l,i}^{\mathsf H}\h_{l,j}|}{\|\h_{l,i}\|\,\|\h_{l,j}\|}\in[0,1].
\end{equation}
Users $i,j$ are \emph{unresolvable at AP $l$} when $\mu_l(i,j)$ is close to one. With a sparse AP this happens (a) when they share the angle and are too close in range (fixed by the larger aperture), or (b) when one sits at the ambiguity point of the other (new, caused by grating lobes). A useful network-level quantity is
\begin{equation}
  \mu_{\rm net}(i,j)=\min_{l\in\mathcal T}\mu_l(i,j),
\end{equation}
the best separation available over a set of APs $\mathcal T$. For $\mathcal T=\mathcal S_k$ (cluster) versus $\mathcal T=\{1,\dots,L\}$ (all APs), the difference is exactly what cross-AP detection exploits.

\subsection{Effective degrees of freedom}
For a channel matrix $\mathbf H$ with Gram matrix $\mathbf R=\mathbf H\mathbf H^{\mathsf H}$,
\begin{equation}
  \mathrm{EDoF}=\frac{\big(\mathrm{tr}\,\mathbf R\big)^2}{\|\mathbf R\|_F^2},
\end{equation}
which equals the number of equally strong eigen-channels. This is the metric used in [5], [8].

\subsection{Effect on your cross-AP information surplus}
Your papers define the per-AP post-combining SINR $\gamma_{l,k}$ and the surplus
\begin{equation}
  \Delta_k=\varepsilon^2_{\mathcal S}-\varepsilon^2_{\mathcal A}\approx
  \frac{\sum_{l\notin\mathcal S_k}\gamma_{l,k}}{\big(1+\sum_{l\in\mathcal S_k}\gamma_{l,k}\big)\big(1+\sum_{l=1}^L\gamma_{l,k}\big)} .
\end{equation}
With a sparse AP, $\gamma_{l,k}$ collapses at APs where user $k$ sits on an ambiguity point of a strong interferer, and is large elsewhere. If the serving cluster happens to contain such an AP, the cluster estimate is unreliable while the all-AP estimate is not: $\Delta_k$ grows. A Transactions-level result would be a bound on $\mathbb E[\Delta_k]$ as a function of $\eta$.

%=====================================================================
\section{Literature map}
\label{sec:lit}

\begin{longtable}{@{}p{0.27\linewidth}p{0.36\linewidth}p{0.30\linewidth}@{}}
\caption{Key works on sparse arrays and near-field cell-free systems.}\\
\toprule
\textbf{Work} & \textbf{What they did} & \textbf{What is missing for you}\\
\midrule
\endfirsthead
\toprule
\textbf{Work} & \textbf{What they did} & \textbf{What is missing for you}\\
\midrule
\endhead
\multicolumn{3}{@{}l}{\textit{Sparse arrays in near-field communications (single array, single cell)}}\\[2pt]
Chen, Ren, Pan, Wang, Wang [5] & Beam analysis of sparse arrays in near-field XL-MIMO; closed-form focusing lobe and grating-lobe suppression; EDoF. Grating lobes are naturally suppressed in the near field, more so for larger spacing and shorter distance. & Single array; no cell-free; no detection or estimation.\\
Chen, Qi, Li, Dobre [6] & Uniform and non-uniform sparse arrays for near-field multi-user; antenna positions optimized for sum rate; channels sparse in a surrogate distance--angle domain, so an OMP estimator is built. & Single array; linear precoding; no SIC or list detection.\\
Zhou, You, Zhang, Chen, Shi [7] & Beam-pattern analysis of sparse arrays in the near field (grating lobes, $2(U-1)$ strong lobes for a linear sparse array) and hybrid beamforming design. & Downlink, single array.\\
Wang, Feng, Zeng, Jin, Yuen, Clerckx, Zhang [8] & Sparse MIMO for near- and far-field users: closed-form near-field EDoF increasing with sparsity; interference suppression. & No cell-free, no estimation.\\
Hussain, Abdallah, Celik, Bj\"ornson, Eltawil [10] & Reconfigurable array thinning for near-field MU-MIMO (PSO), comparable to movable antennas. & Single array.\\
Li, Min, Zeng, Jin, Dai, Yuan, Zhang [9] & Survey: sparse MIMO for ISAC; co-array with $\mathcal O(M^2)$ virtual elements; grating-lobe suppression, codebooks. & Sensing focus.\\
{[15], [16]} (authors not verified) & Sparse arrays for near-field sensing and communications; constant-distance focusing with reduced focal shift. & Single array.\\[4pt]
\multicolumn{3}{@{}l}{\textit{Modular and quasi-distributed arrays (sparse at sub-array level)}}\\[2pt]
Li, Dong, Zeng, Jin, Zhang [11] & Modular XL-array: better resolution than a collocated array with the same elements, but grating lobes; near-field model suppresses them better than the far-field model. & Single modular array.\\
{[12]} (authors not verified) & Multi-user modular XL-MIMO: grating lobes degrade SINR; user grouping (scheduling) to avoid them. & No detection design.\\
Kosasih, Demir, Bj\"ornson [13] & Modular linear arrays: ideal near-field focusing by choosing sub-array sizes; localization by fusing per-ULA estimates; parametric channel estimation. & Not cell-free, but the ``fuse per-sub-array estimates'' idea is close to your cross-AP fusion.\\
{[14]} (authors not verified) & Grating lobes for quasi-distributed arrays, from ULA to MRA configurations. & Single array.\\[4pt]
\multicolumn{3}{@{}l}{\textit{Near-field cell-free (no sparse-array focus, except [25])}}\\[2pt]
Wang \emph{et al.} [21] & Distributed combining for near-field cell-free XL-MIMO (MMSE/L-MMSE, SSOR). & Dense ULAs; linear combining only.\\
Lei \emph{et al.} [22] & Distributed activity detection in hybrid near-far field cell-free. & Activity detection, not data detection.\\
Lei, Wang, Xiao, Zhang, Ai [23] & Uplink performance of cell-free XL-MIMO. & Dense arrays.\\
Mylonopoulos, Interdonato, Buzzi, Liu [24] & Near-field beamfocusing in cell-free networks exploiting LoS from distributed APs. & Downlink focusing.\\
Xie \emph{et al.} [25] & Beam focusing for near-field cell-free massive MIMO with few antennas per AP; delay-based focusing (OFDM) plus a sparse-array angle-based focusing scheme. & The only near-field cell-free work with sparse arrays found; downlink focusing, no detection, no estimation.\\[4pt]
\multicolumn{3}{@{}l}{\textit{Foundations you already cite or should cite}}\\[2pt]
Cui, Dai [1]; Wei, Dai [33] & Polar-domain dictionary and OMP (P-SOMP, P-SIGW); hybrid-field estimation. & Dense arrays; dictionary must be redesigned for sparse spacing.\\
Lu, Zeng \emph{et al.} [2]; Wang \emph{et al.} [30] & Tutorials on near-field XL-MIMO; sparse and modular arrays covered. & ---\\
Moffet [18]; Pal, Vaidyanathan [19]; Vaidyanathan, Pal [20] & MRA, nested and co-prime arrays. & Classic array processing.\\
\bottomrule
\end{longtable}

%=====================================================================
\section{The gap}
\label{sec:gap}

Within the limits of this search (see the box on page~1):
\begin{enumerate}
  \item Sparse arrays in the near field are studied for \textbf{single arrays}: beam patterns, EDoF, precoding, thinning, sensing [5]--[11], [13].
  \item Near-field \textbf{cell-free} work uses dense arrays and focuses on combining, activity detection, or downlink focusing [21]--[24]. Only [25] uses a sparse array, and only for downlink beam focusing.
  \item \textbf{No work was found} on \emph{uplink data detection} (SIC, list detection, IDD) or on \emph{channel estimation} for \emph{hybrid-field cell-free} networks with \emph{sparse APs}.
  \item No work was found that analyses how the per-link near/far-field classification \emph{and} the per-AP grating-lobe ambiguity interact across APs, nor one that uses \emph{multi-AP fusion} to resolve grating-lobe ambiguities.
\end{enumerate}
Point 4 is the strongest claim you can make, because it connects directly to the cross-AP idea of your two conference papers.

%=====================================================================
\section{Options for bringing sparse arrays into your model}
\label{sec:options}

\begin{longtable}{@{}p{0.2\linewidth}p{0.4\linewidth}p{0.33\linewidth}@{}}
\toprule
\textbf{Option} & \textbf{What changes} & \textbf{Pros and cons}\\
\midrule
\endhead
A. Uniform sparse ULA, same $N$ & Spacing $\eta\lambda/2$ at every AP; $R$ grows by $\eta^2$; range resolution improves (Table~\ref{tab:numbers}). & + Easiest; biggest gain in range separation. $-$ Larger AP; hybrid regime may vanish (rule R2); grating lobes.\\
B. Thinned, same aperture & $N/\eta$ antennas over the same 3.15\,m. & + Same $R$, same hybrid regime; $\eta\times$ fewer antennas, cheaper combiners ($\mathcal O(N^3)\to\mathcal O(N^3/\eta^3)$). $-$ Grating lobes with no resolution gain; lower array gain.\\
C. Non-uniform (nested, co-prime, MRA) & Positions from [18]--[20]. & + No exact grating lobes; larger co-array helps estimation. $-$ Higher side-lobes; harder analysis.\\
D. Modular AP & Dense sub-arrays with gaps [11], [13]. & + Practical to build; known theory. $-$ Grating lobes at sub-array level.\\
E. Heterogeneous APs & Some APs dense, some sparse (or $\eta$ chosen per AP). & + Keeps a true hybrid-field network; lets you study which APs should be sparse. $-$ Larger design space.\\
F. Network as a sparse array & Treat the $L$ APs as one very sparse distributed array (positions $\mathbf r_l$). & + This is already what cross-AP fusion exploits; gives a clean theory story. $-$ Needs care with synchronisation assumptions.\\
\bottomrule
\end{longtable}

\begin{keybox}{Recommendation}
Use \textbf{Option A} as the main study ($\eta\in\{1,2,4\}$, same $N$), and \textbf{Option B} as the fairness check (same aperture). Add \textbf{Option E} for a short ``design insight'' figure if space allows. Option C is a good follow-up paper.
\end{keybox}

%=====================================================================
\section{Grating-lobe-aware channel estimation (connects to your VTC paper)}
\label{sec:est}

Your VTC estimator (polar-domain OMP) must change in three ways for a sparse AP:
\begin{enumerate}
  \item \textbf{Finer angular grid.} Beams are $\eta$ times narrower, so use about $\eta\cdot o\cdot N$ angle samples ($o$ = oversampling).
  \item \textbf{Alias testing.} Atoms at $s$ and $s\pm 2/\eta$ share the same linear phase and differ only through the curvature term. After OMP picks an atom $(s,r)$, also refine the alias candidates $(s\pm 2m/\eta,\, r')$ and keep the best fit. (The current authentic estimator in your MATLAB code already does a simple version of this at end-fire, where $s=-1$ and $s=+1$ are the same direction for $\eta=1$.)
  \item \textbf{Cross-AP consistency (new, cell-free specific).} A user has \emph{one} position $\mathbf u_k$, and it determines $(\theta_{l,k},d_{l,k})$ at every AP. Instead of estimating a polar atom per AP independently, estimate $\mathbf u_k$ jointly from all near-field APs: an alias that fits AP $l$ will not be consistent with the position seen by AP $l'$. This is a \emph{location-domain dictionary}, related to the per-sub-array fusion of [13], and it is new for cell-free networks.
\end{enumerate}
A clean analytical result to aim for: the coherence of the polar dictionary as a function of $\eta$, and the NMSE gain of the location-domain estimator over per-AP OMP.

%=====================================================================
\section{A plan for the Transactions paper}
\label{sec:plan}

\paragraph{Working title.} \emph{Sparse-Array Near-Field Cell-Free XL-MIMO: Grating-Lobe-Aware Channel Estimation and Cross-AP List Detection with Iterative Decoding.}

\paragraph{Contributions (draft).}
\begin{enumerate}[label=\textbf{C\arabic*.}]
  \item Hybrid-field cell-free model with sparse APs; closed-form fraction of near-field links and of finite-depth links as a function of $\eta$ and the area.
  \item Near-field grating-lobe analysis for cell-free: ambiguity points per AP, and the probability that a user pair is unresolvable at one AP, in a cluster, and at all APs.
  \item Grating-lobe-aware estimation: aliased polar dictionary plus cross-AP location-domain refinement; NMSE analysis.
  \item Cross-AP list SIC and IDD (from ICASSP and VTC), with the information surplus $\Delta_k$ analysed as a function of $\eta$.
  \item Clustering that accounts for grating-lobe interference (see the separate clustering document).
\end{enumerate}

\paragraph{Figures to plan.}
\begin{enumerate}
  \item Fraction of near-field and finite-depth links vs.\ $\eta$ (analysis and simulation).
  \item Beam pattern of one AP, $\eta=1,2,4$: main lobe, grating lobe at the same range, ambiguity point.
  \item NMSE vs.\ SNR: per-AP OMP, aliased OMP, location-domain estimator; $\eta=1,2,4$.
  \item BER vs.\ SNR for the four detectors, $\eta=1,2,4$, perfect and estimated CSI (extension of VTC Fig.~1).
  \item Same-$N$ vs.\ same-aperture comparison (Options A and B).
  \item $\Delta_k$ distribution (CDF) vs.\ $\eta$; gate firing rate and fronthaul vs.\ SNR.
  \item IDD BER at iteration 3, $\eta=1,2,4$.
\end{enumerate}

\paragraph{Code changes in your MATLAB script (step by step).}
\begin{enumerate}
  \item Add \texttt{eta} near the top; set \texttt{d\_ant = eta*lambda/2} and \texttt{d\_H = eta/2} (the far-field correlation uses spacing in wavelengths). Keep $N$ (Option A) or set \texttt{N = round((N0-1)/eta)+1} (Option B).
  \item $D_{\rm AP}$, \texttt{d\_Ray} and the near-field mask then update automatically; print the near-field fraction.
  \item Build the polar dictionary with \texttt{pd\_os\_ang} scaled by \texttt{eta}; add alias candidates $s\pm 2m/\eta$ to the refinement starts in \texttt{pdomp\_estimate}; change the wrap period in \texttt{refine\_param} accordingly.
  \item Apply the same spacing to the centralized-BS benchmark, or state clearly that it stays dense.
  \item Keep \texttt{genie\_csi = false}, so that every estimated-CSI result is authentic.
\end{enumerate}

\paragraph{Risks and honest notes.}
\begin{itemize}
  \item With $\eta\ge2$ in a $300\times300$\,m$^2$ area nearly all links are near-field, so ``hybrid-field'' must be argued with a larger area or with mixed APs.
  \item A larger aperture also increases spatial non-stationarity (visibility regions); your model assumes every element sees the user, so state it.
  \item Sparse arrays help range separation only for users inside $R/\alpha$; report the fraction of such users.
\end{itemize}

%=====================================================================
\begin{thebibliography}{99}\small
\bibitem{cui} M.~Cui and L.~Dai, ``Channel estimation for extremely large-scale MIMO: Far-field or near-field?'' \emph{IEEE Trans.\ Commun.}, vol.~70, no.~4, pp.~2663--2677, 2022.
\bibitem{lu} H.~Lu, Y.~Zeng, C.~You, Y.~Han, J.~Zhang, Z.~Wang, Z.~Dong, S.~Jin, C.-X.~Wang, T.~Jiang \emph{et al.}, ``A tutorial on near-field XL-MIMO communications toward 6G,'' \emph{IEEE Commun.\ Surveys Tuts.}, vol.~26, no.~4, pp.~2213--2257, 2024.
\bibitem{primer} E.~Bj\"ornson, \"O.~T.~Demir, and L.~Sanguinetti, ``A primer on near-field beamforming for arrays and reconfigurable intelligent surfaces,'' in \emph{Proc.\ Asilomar Conf.}, 2021; arXiv:2110.06661.
\bibitem{chapter} ``Near-field beamforming and multiplexing using extremely large aperture arrays,'' book chapter, arXiv:2209.03082, 2022 (authors not verified).
\bibitem{chenren} X.~Chen, H.~Ren, C.~Pan, C.-X.~Wang, and J.~Wang, ``Exploring the advantages of sparse arrays in near-field XL-MIMO systems: Beam analysis and EDoF function,'' arXiv:2501.09234, 2025.
\bibitem{chenqi} K.~Chen, C.~Qi, G.~Y.~Li, and O.~A.~Dobre, ``Near-field multiuser communications based on sparse arrays,'' \emph{IEEE J.\ Sel.\ Topics Signal Process.}, vol.~18, no.~4, 2024; arXiv:2406.09238.
\bibitem{zhou} C.~Zhou, C.~You, H.~Zhang, L.~Chen, and S.~Shi, ``Sparse array enabled near-field communications: Beam pattern analysis and hybrid beamforming design,'' \emph{IEEE Trans.\ Wireless Commun.}, 2025, doi:10.1109/TWC.2025.3578561; arXiv:2401.05690.
\bibitem{wangfeng} H.~Wang, C.~Feng, Y.~Zeng, S.~Jin, C.~Yuen, B.~Clerckx, and R.~Zhang, ``Enhancing spatial multiplexing and interference suppression for near- and far-field communications with sparse MIMO,'' arXiv:2408.01956, 2024.
\bibitem{lisparse} X.~Li, H.~Min, Y.~Zeng, S.~Jin, L.~Dai, Y.~Yuan, and R.~Zhang, ``Sparse MIMO for ISAC: New opportunities and challenges,'' arXiv:2406.12270, 2024.
\bibitem{hussain} A.~Hussain, A.~Abdallah, A.~Celik, E.~Bj\"ornson, and A.~M.~Eltawil, ``Sparse array design for near-field MU-MIMO: Reconfigurable array thinning approach,'' arXiv:2602.21973, 2026 (reported as accepted in IEEE Wireless Commun.\ Lett.).
\bibitem{limod} X.~Li, Z.~Dong, Y.~Zeng, S.~Jin, and R.~Zhang, ``Near-field beam focusing pattern and grating lobe characterization for modular XL-array,'' arXiv:2305.05408, 2023.
\bibitem{mumod} ``Multi-user modular XL-MIMO communications: Near-field beam focusing pattern and user grouping,'' arXiv:2308.11289, 2023 (authors not verified).
\bibitem{kosasih} A.~Kosasih, \"O.~T.~Demir, and E.~Bj\"ornson, ``Near-field beamfocusing, localization, and channel estimation with modular linear arrays,'' arXiv:2505.07991, 2025.
\bibitem{quasi} ``Near-field communications with grating lobes for quasi-distributed arrays: From ULA to MRA,'' arXiv:2607.29341, 2026 (authors not verified).
\bibitem{enh} ``Enhancing near-field sensing and communications with sparse arrays: Potentials, challenges, and emerging trends,'' arXiv:2309.08681, 2023 (authors not verified).
\bibitem{const} ``Sparse arrays enable near-field constant-distance focusing with reduced focal shift,'' arXiv:2505.07285, 2025 (authors not verified).
\bibitem{coprimeloc} ``Near-field localization with coprime array,'' arXiv:2411.01529, 2024 (authors not verified).
\bibitem{moffet} A.~T.~Moffet, ``Minimum-redundancy linear arrays,'' \emph{IEEE Trans.\ Antennas Propag.}, vol.~16, no.~2, pp.~172--175, 1968.
\bibitem{nested} P.~Pal and P.~P.~Vaidyanathan, ``Nested arrays: A novel approach to array processing with enhanced degrees of freedom,'' \emph{IEEE Trans.\ Signal Process.}, vol.~58, no.~8, pp.~4167--4181, 2010.
\bibitem{coprime} P.~P.~Vaidyanathan and P.~Pal, ``Sparse sensing with co-prime samplers and arrays,'' \emph{IEEE Trans.\ Signal Process.}, vol.~59, no.~2, pp.~573--586, 2011.
\bibitem{wangcomb} Z.~Wang, J.~Zhang, B.~Xu, D.~Niyato, B.~Ai, S.~Mao, and Z.~Han, ``Low-complexity distributed combining design for near-field cell-free XL-MIMO systems,'' \emph{IEEE Trans.\ Wireless Commun.}, vol.~25, pp.~11799--11815, 2026; arXiv:2602.03581.
\bibitem{lei} J.~Lei, Y.~Li, Z.~Wang, Q.~Lin, Y.-F.~Liu, and Y.-C.~Wu, ``A unified distributed algorithm for hybrid near-far field activity detection in cell-free massive MIMO,'' \emph{IEEE Trans.\ Wireless Commun.}, 2026; arXiv:2509.15162.
\bibitem{leiul} H.~Lei, Z.~Wang, H.~Xiao, J.~Zhang, and B.~Ai, ``Uplink performance of cell-free extremely large-scale MIMO systems,'' arXiv:2301.12086, 2023.
\bibitem{mylo} G.~Mylonopoulos, G.~Interdonato, S.~Buzzi, and P.~Liu, ``Leveraging line-of-sight propagation for near-field beamfocusing in cell-free networks,'' arXiv:2503.21599, 2025.
\bibitem{xie} M.~Xie \emph{et al.}, ``Beam focusing for near-field cell-free massive MIMO,'' \emph{IEEE Wireless Commun.\ Lett.}, 2025, doi:10.1109/LWC.2025.3574325 (full author list not verified).
\bibitem{bjsc} E.~Bj\"ornson and L.~Sanguinetti, ``Scalable cell-free massive MIMO systems,'' \emph{IEEE Trans.\ Commun.}, vol.~68, no.~7, pp.~4247--4261, 2020.
\bibitem{found} \"O.~T.~Demir, E.~Bj\"ornson, and L.~Sanguinetti, ``Foundations of user-centric cell-free massive MIMO,'' \emph{Found.\ Trends Signal Process.}, vol.~14, no.~3--4, pp.~162--472, 2021.
\bibitem{ngo} H.~Q.~Ngo, A.~Ashikhmin, H.~Yang, E.~G.~Larsson, and T.~L.~Marzetta, ``Cell-free massive MIMO versus small cells,'' \emph{IEEE Trans.\ Wireless Commun.}, vol.~16, no.~3, pp.~1834--1850, 2017.
\bibitem{haider} N.~Haider and R.~C.~de~Lamare, ``Iterative interference cancellation for mixed near- and far-field XL-MIMO systems,'' \emph{IEEE Wireless Commun.\ Lett.}, vol.~15, pp.~5323--5327, 2026.
\bibitem{wangtut} Z.~Wang, J.~Zhang, H.~Du, D.~Niyato, S.~Cui, B.~Ai, M.~Debbah, K.~B.~Letaief, and H.~V.~Poor, ``A tutorial on extremely large-scale MIMO for 6G: Fundamentals, signal processing, and applications,'' \emph{IEEE Commun.\ Surveys Tuts.}, 2024.
\bibitem{stutz} A.~Stutz-Tirri, A.~Dkhan, H.~Sarieddeen, and C.~Studer, ``On near-far-field boundaries in wireless systems,'' arXiv:2605.02424, 2026.
\bibitem{zhu} L.~Zhu, W.~Ma, and R.~Zhang, ``Movable antennas for wireless communication: Opportunities and challenges,'' \emph{IEEE Commun.\ Mag.}, vol.~62, no.~6, pp.~114--120, 2024.
\bibitem{wei} X.~Wei and L.~Dai, ``Channel estimation for extremely large-scale massive MIMO: Far-field, near-field, or hybrid-field?'' \emph{IEEE Commun.\ Lett.}, vol.~26, no.~1, pp.~177--181, 2022.
\end{thebibliography}

\end{document}
